%!TEX root = ../main.tex

\section{Адаптация расчетного шага по времени}

\subsection{Общий вид формулы адаптивного шага}

Одна из проблем, возникающих при работе с моделью, состоит в следующем. Характерным поведением системы является резкое, стремительное образование канала пробоя, которому предшествует длительный период очень слабых и медленных изменений в системе. Шаг расчетной сетки по времени должен быть достаточно мал для адекватного моделирования быстрых процессов (косвенно это выражается, например, условиями устойчивости схемы, сформулированными ранее), однако настолько сильное разрешение по времени оказывается избыточным в период медленной эволюции системы. Таким образом, использование регулярной по времени сетки видится нерациональным.

В этом разделе будет исследовано несколько методов адаптации расчетного шага по времени в применении к разностной схеме \eqref{sch:one_dim} для задачи \eqref{eq:one_dim} в одномерной, упрощенной, постановке. Однако существенно, что методы не опираются на размерность и в дальнейшем будут применены для ускорения расчета задачи в полной постановке. Везде в этом разделе считается заданным условие \eqref{cond:boundary_voltage}, то есть $E = \Const$.

При переходе между временными слоями $k$ и $k + 1$ будем использовать переменный шаг по времени $\tau^k$. Уравнение \eqref{sch:one_dim_a} приобретает вид
\begin{equation}
	\phi_j^{k + 1} = \phi_j^k + m \tau^k \left[ \half (E^k)^2 \epsilon'(\phi_j^k) + \cfrac{\Gamma}{l^2} f'(\phi_j^k) + \cfrac{\Gamma}{2} \cfrac{\phi_{j + 1}^k - 2 \phi_j^k + \phi_{j - 1}^k}{h^2} \right],
	\label{sch:transition_adaptive}
\end{equation}

Далее будет рассмотрено несколько подходов к вычислению $\tau^k$. Общая формула его расчета
\begin{equation}
	\tau^k = \max \left[ \taumin, \min( \taumax, \widetilde{\tau}^k) \right],
	\label{sch:time_step_min_max}
\end{equation}
то есть величина $\widetilde{\tau}^k$ (рассчитываемая по своей формуле для каждого подхода) ограничивается снизу и сверху заранее выбранными значениями $\taumin$ и $\taumax$ соответственно. Последние выбираются эмпирически, исходя из критериев устойчивости, точности и скорости расчета.


\subsection{Адаптации по фазовому полю и по энергии}

Рассмотрим первые два подхода к адаптации шага по времени, предложенные в статьях \cite{li_time_step} и \cite{zhang_time_step} соответственно. Введем их вместе из-за определенной общности:
\begin{align}
	\widetilde{\tau}_1^k & = \cfrac{\tol_1}{\left\| \left[ \partflt{\phi} \right]_h \right\|_C},
	\label{sch:time_step_phi} \\
	\widetilde{\tau}_2^k & = \cfrac{\tol_2}{\left| \left[ d \Psi / dt \right]_h \right|}.
	\label{sch:time_step_energy}
\end{align}
Числовые константы $\tol_1$ и $\tol_2$ подбираются эмпирически. Здесь и далее символами $[\, \cdot \,]_h$ и $[\, \cdot \,]_j^k$ обозначаются сеточные функции, в большинстве случаев~-- разностные производные ($j$, $k$ -- индексы на сетке, целые либо полуцелые).

В формуле \eqref{sch:time_step_phi} в качестве $[\partflt{\phi}]_h$ удобно использовать
\[
	\left[ \partt{\phi} \right]_j^{k + 1/2} = \cfrac{\phi_j^{k + 1} - \phi_j^k}{\tau^k}
\]
из левой части разностного уравнения \eqref{sch:one_dim_task_a}. Дополнительные вычисления не потребуются, так как в расчете уже используется значение правой части уравнения~\eqref{sch:one_dim_task_a}. Если такой подход неприменим (например, из-за трудностей с синхронизацией параллельных вычислений), то можно использовать производную $[\partflt{\phi}]^{k - 1/2}$, сохраненную с предыдущего временного слоя.

В формуле \eqref{sch:time_step_energy} в знаменателе находится модуль разностной производной свободной энергии $\Psi(t)$ (см. \eqref{eq:energy_density_free}). Согласно замечанию \ref{rem:Phi_constant} в адекватном расчете рассматриваемой постановки задачи $[d \Psi / dt]_h$ либо отрицательна, либо крайне мала по модулю (сеточный эффект колебания системы вблизи минимума $\Psi$).

Для расчета плотности $\psi$ свободной энергии необходима разностная производная $[\partflx{\phi}]_h$ (см. выражения \eqref{eq:energy_densities}, $\beta = 0$). Предлагается использовать следующие формулы:
\begin{gather*}
	\allowdisplaybreaks
	\psi_j^k = F(\phi_j^k) + \cfrac{\Gamma}{4} \left( \left[ \partx{\phi} \right]_j^k \right)^2, \qquad \left[ \partx{\phi} \right]_j^k = \begin{cases}
		\cfrac{\phi_1^k - \phi_0^k}{h} & \text{при } j = 0; \\
		\cfrac{\phi_{j + 1}^k - \phi_{j - 1}^k}{2 h} & \text{при } j = \overline{1, M - 1}; \\
		\cfrac{\phi_M^k - \phi_{M - 1}^k}{h} & \text{при } j = M;
	\end{cases} \\
	\Psi^k = \cfrac{h \psi_0^k + h \psi_M^k}{2} + \sum\limits_{j = 1}^{M - 1} h \psi_j^k = h \cdot \left( \cfrac{\psi_0^k + \psi_M^k}{2} + \sum\limits_{j = 1}^{M - 1} \psi_j^k \right).
\end{gather*}
В формуле \eqref{sch:time_step_energy} будем использовать разностную производную энергии с предыдущего временного слоя
\[
	\left[ \cfrac{d \Psi}{dt} \right]^{k - 1/2} = \cfrac{\Psi^k - \Psi^{k - 1}}{\tau^{k - 1}},
\]
положив $\tau^0 = \taumin$.

Использование формулы \eqref{sch:time_step_phi} для расчета $\widetilde{\tau}^k$ будем условно называть адаптацией по фазовому полю, формулы \eqref{sch:time_step_energy} -- адаптацией по энергии.

Общность описанных двух подходов заключается в следующем. Для адаптации по фазовому полю рассмотрим норму приращения $\phi^{k + 1} - \phi^k$:
\[
	\left\| \phi^{k + 1} - \phi^k \right\|_C = \tau_1^k \cdot \left\| \left[ \partt{\phi} \right]^{k + 1/2} \right\|_C = \tol_1 \left\| \left[ \partt{\phi} \right]^{k + 1/2} \right\|_C^{-1} \cdot \left\| \left[ \partt{\phi} \right]^{k + 1/2} \right\|_C = \tol_1.
\]
Выходит нормирование приращения фазового поля (с оговоркой на ограничения через $\taumin$ и $\taumax$).

В случае адаптации по энергии можно провести похожее рассуждение. В рассматриваемой постановке задачи выполняется равенство \eqref{eq:energy_free_decrease}, что для сеточных функций дает
\[
	\left| \left[ \cfrac{d \Psi}{dt} \right]_h \right| \approx \cfrac{1}{m} \left\| \left[ \partt{\phi} \right]_h \right\|_2^2.
\]
Здесь $\| \cdot \|_2$ обозначена сеточная $L_2$-норма:
\[
	\| [\zeta]_h \|_2 = \sqrt{h \sum\limits_{j = 0}^M (\zeta_j)^2}.
\]
Таким образом, при адаптации по~энергии
\begin{multline*}
	\left\| \phi^{k + 1} - \phi^k \right\|_2 = \tau_2^k \cdot \left\| \left[ \partt{\phi} \right]^{k + 1/2} \right\|_2 \approx \tol_2 \cdot m \left\| \left[ \partt{\phi} \right]^{k - 1/2} \right\|_2^{-2} \cdot \left\| \left[ \partt{\phi} \right]^{k + 1/2} \right\|_2 \approx \\ \approx \tol_2 \cdot m \cdot \left\| \left[ \partt{\phi} \right]_h \right\|_2^{-1}.
\end{multline*}

Авторы не стали отклоняться от предложенного в работе \cite{zhang_time_step} метода адаптации, однако из проделанных рассуждений получается, что из модуля производной энергии в формуле \eqref{sch:time_step_energy} логичнее было бы извлечь квадратный корень, чтобы выполнялось $\| \phi^{k + 1} \hm - \phi^k \|_2 \approx \tol_2 \cdot \sqrt{m}$.